\section{Chapter 3: Web Application}\label{chapter-3-web-application}

\subsection{3.1 Requirements Analysis}\label{requirements-analysis}

The system requirements are defined in \texttt{docs/SRS.md}, focusing on two primary actors: the \textbf{User} and the \textbf{Data Owner}.

\subsubsection{3.1.1 Functional Requirements}\label{functional-requirements}

\begin{itemize}
\tightlist
\item
  \textbf{UC-RETRIEVE-01}: The core workflow where users upload images, review segmentation, and receive ranked species predictions.
\item
  \textbf{UC-DATA-01}: Allows the Data Owner to index new reference data, ensuring the retrieval database stays current.
\item
  \textbf{UC-MODEL-01}: Provides tools for the Data Owner to evaluate and promote ``Candidate Models'' to the production index.
\end{itemize}

\subsubsection{3.1.2 Non-Functional Requirements}\label{non-functional-requirements}

\begin{itemize}
\tightlist
\item
  \textbf{Performance}: Retrieval must complete within 5 seconds.
\item
  \textbf{Security}: Implementation of JWT-based authentication and Role-Based Access Control (RBAC) to protect sensitive dataset mutations.
\item
  \textbf{Auditability}: Every change made by a Data Owner is recorded in an audit log.
\end{itemize}

\subsection{3.2 High-Level Architecture}\label{high-level-architecture}

The system follows a decoupled client-server architecture: - \textbf{Frontend}: Built with \textbf{React 19} and \textbf{Vite}, utilizing a scientist-facing design language for high-density data display. - \textbf{Backend}: A \textbf{FastAPI} server that orchestrates communication between the frontend, the vector store, and the relational database.

\subsection{3.3 Database Design \& Optimization}\label{database-design-optimization}

\subsubsection{3.3.1 Relational Database (PostgreSQL)}\label{relational-database-postgresql}

Used for structured data: - \texttt{users}: Authentication and role management. - \texttt{species} \& \texttt{media}: Managed catalogs of fungal taxa and growth media. - \texttt{dataset\_items}: Metadata linking images to their species and strains.

\subsubsection{3.3.2 Vector Database (Qdrant)}\label{vector-database-qdrant}

Used for high-dimensional search: - Stores embeddings of the reference dataset. - Enables payload filtering (e.g., \texttt{where\ media\ ==\ \textquotesingle{}PDA\textquotesingle{}}) during the KNN search.

\subsubsection{3.3.3 Why Dual-DB?}\label{why-dual-db}

A single database cannot efficiently handle both complex relational queries and high-dimensional vector searches. PostgreSQL provides the necessary consistency for governance, while Qdrant provides the sub-second latency required for retrieval.

\subsection{3.4 API Design}\label{api-design}

The API utilizes an asynchronous job pattern for long-running retrieval tasks to prevent request timeouts.

\begin{lstlisting}[language=Python, caption=Asynchronous Retrieval Endpoint]
@router.post("/query", response_model=RetrievalJobResponse, status_code=202)
def start_query(data: RetrievalQueryRequest, user: User = Depends(get_current_user)) -> dict:
    job_id = new_id()
    job = {
        "id": job_id,
        "status": "processing",
        "config": data.model_dump(),
        "created_at": utcnow(),
    }
    get_retrieval_job_store().put(job)
    return {"job_id": job_id, "status": "processing"}
\end{lstlisting}

\subsubsection{3.4.1 Authentication}\label{authentication}

JWT implementation and RBAC.

\subsubsection{3.4.2 Retrieval Endpoint}\label{retrieval-endpoint}

Design of the \texttt{/retrieve} endpoint.

\subsubsection{3.4.3 Indexing Endpoint}\label{indexing-endpoint}

Design of the \texttt{/index} endpoint.

\subsection{3.5 UI/UX Design}\label{uiux-design}

The UI is optimized for a research workflow: - \textbf{Upload Interface}: Supports both single-image and batch-folder uploads. - \textbf{Segmentation Review}: An interactive canvas allowing users to drag and resize bounding boxes before final prediction. - \textbf{Results View}: Displays ranked species with confidence scores and visual matches from the reference set.
